Telecommunications scam and phishing campaigns in Vietnam inflict severe damage by exploiting localized syntactic nuances like teencode and non-diacritic text. 
Current defensive solutions either depend on symmetric-loss classifiers suffering from unacceptable false-negative rates, or rely on cloud-based Large Language Models (LLMs) exhibiting severe alarm fatigue and latency bottlenecks. 
In this work, we conducted an empirical study benchmarking pretrained Vietnamese language models against in-context LLMs, and proposed a lightweight edge screener that pairs ViSoBERT with a Cost-Sensitive Weighted Binary Cross-Entropy loss ($\mathcal{L}_{\text{WBCE}}$, $\alpha=5.0$) quantized via ONNX INT8. Evaluated across $5$ independent random seeds on an isolated frozen test partition ($N=267$) from a corpus of $2,664$ real-world Vietnamese SMS messages, the proposed ViSoBERT model achieved a Scam Recall of $95.62\% \pm 2.25\%$ and a Macro-F1 score of $92.51\% \pm 2.39\%$. 
It significantly outperformed the Gemini 3.5 Flash baseline ($p < 0.0001$, McNemar test), whose precision collapsed to $64.47\%$ due to hyper-vigilant over-flagging, while maintaining an on-device CPU execution latency of $15.8\text{ ms/message}$ and a $390\text{ MB}$ footprint. 
These results demonstrate that asymmetric penalty weighting on social-specialized encoders provides Pareto-optimal edge triage, establishing the empirical grounding for a two-tier multi-agent hybrid defense architecture.
